\begin{enumerate}\small
\item Toute personne qui, volontairement, porte des coups à autrui est punie d'un emprisonnement d'un an et d'une amende de mille dinars. \hfill{\footnotesize[PenalFr/s02/218]}
\item Si l'auteur des coups est un descendant ou le conjoint de la victime, la peine est de deux ans d'emprisonnement et de deux mille dinars d'amende. \hfill{\footnotesize[PenalFr/s02/218]}
\item En cas de préméditation, la peine est de trois ans d'emprisonnement et de trois mille dinars d'amende. \hfill{\footnotesize[PenalFr/s02/218]}
\item La tentative d'administration de coups est également punissable. \hfill{\footnotesize[PenalFr/s02/218]}
\item Toute personne ayant participé à une rixe au cours de laquelle des coups ont été portés encourt un emprisonnement de six mois, sans préjudice des peines prévues à l'article 218 contre l'auteur des coups. \hfill{\footnotesize[PenalFr/s02/220]}
\end{enumerate}
